\chapter{Mental Arithmetic}\label{ch:iv:mental-arithmetic}

Across the desk: ``Seven per cent of 340 million, spread over 250 trading days, per day?'' The answer, about
95 thousand, comes back in five seconds from the candidate who turned seven per cent over 250 days into 2.8
basis points a day, and 2.8 basis points of 340 million into 340 times 280; the one who started long
multiplication in her head is still carrying digits. Speed in arithmetic is not a talent for carrying digits;
it is a small set of rewrites that turn a hard calculation into an easy one, practised until they are
automatic, and a habit of checking the result.

\section{Products and squares}

Most fast products use one of four rewrites. Each replaces a multiplication by one that is done in a step.

\begin{method}[Four rewrites for products]\label{met:iv:mental-arithmetic:products}
\begin{enumerate}
\item \emph{Difference of squares}: $(a+d)(a-d) = a^2 - d^2$ when the factors are equally spaced around a round
number: $43 \times 37 = 40^2 - 3^2 = 1\,591$.
\item \emph{Near a hundred}: $(100+a)(100+b) = 100(100 + a + b) + ab$, valid for negative $a, b$ too: $103
\times 108 = 11\,100 + 24 = 11\,124$.
\item \emph{Squares ending in 5}: $(10n + 5)^2 = 100\,n(n+1) + 25$: $85^2 = 7\,200 + 25 = 7\,225$.
\item \emph{Friendly factors}: multiply by 25 as by 100 then divide by 4, by 125 as by 1000 then divide by 8,
by 5 as by 10 then halve: $48 \times 25 = 4\,800/4 = 1\,200$.
\end{enumerate}
\end{method}

A product of two numbers that fit none of these is split: $67 \times 38 = 67 \times 40 - 67 \times 2 = 2\,680 -
134 = 2\,546$. The choice of split is the skill: subtract from the nearest round number rather than add
partial products.

\section{Fractions, percentages and reciprocals}

Desk arithmetic is mostly percentages of large numbers and conversions between rates. Three facts make it
fast. A percentage is symmetric: $a\%$ of $b$ equals $b\%$ of $a$, so 18\% of 50 is 50\% of 18, which is 9.
A basis point is $10^{-4}$, so a rate in basis points times a notional in millions is a number of hundreds of
currency units: 2.8 basis points of 340 million is $2.8 \times 340 \times 100 = 95\,200$. And a small table of
reciprocals turns every division by a small integer into a multiplication.

\begin{center}
\small
\begin{tabular}{@{}cccccccccc@{}}
\toprule
$n$ & 6 & 7 & 8 & 9 & 11 & 12 & 13 & 16 & 17 \\
$1/n$ & 0.1667 & 0.1429 & 0.125 & 0.1111 & 0.0909 & 0.0833 & 0.0769 & 0.0625 & 0.0588 \\
\bottomrule
\end{tabular}

\smallskip
{\small Reciprocals worth knowing; $1/7 = 0.\overline{142857}$ repeats with period six and $1/19 \approx
0.0526$ is the next one to learn.}
\end{center}

\section{Roots, logarithms and compounding}

A square root near a known square is linear to first order: $\sqrt{a^2 + d} \approx a + d/(2a)$, with an error
below $d^2/(8a^3)$ for $d \ge 0$. $\sqrt{103} \approx 10 + 3/20 = 10.15$, against the true 10.1489. For
logarithms, three values and one series cover most questions: $\ln 2 \approx 0.693$, $\ln 3 \approx 1.099$,
$\ln 10 \approx 2.303$, and $\ln(1 + x) = x - x^2/2 + x^3/3 - \dots$ for small $x$.

\begin{definition}[Rule of 72]\label{def:iv:mental-arithmetic:rule72}
The \emph{rule of 72}\index{rule of 72} approximates the number of periods an amount takes to double at a rate
of $r$ per cent a period, compounded once a period, by $72/r$. The exact doubling time is $\ln 2 / \ln(1 +
r/100)$.
\end{definition}

\begin{proposition}[Where the rule of 72 is exact]\label{prop:iv:mental-arithmetic:rule72}
The \omterm{def:iv:mental-arithmetic:rule72}{rule of 72} is exact at a rate of about 7.85\%, overestimates the doubling time below it and
underestimates above it. The rule of 70 is exact near 2\%. The rule of 69.3 is exact for continuously
compounded rates ($100\ln 2 \approx 69.3$) and always underestimates for rates compounded once a period.
\end{proposition}

\begin{proof}
$\ln(1 + x) < x$ for $x > 0$, so $\ln 2/\ln(1+x) > \ln 2/x$, which is the last statement. The numerator 72
exceeds $100 \ln 2$ to compensate for $\ln(1+x) < x$; the two effects balance where $72/(100x) = \ln 2/\ln(1+x)$,
solved numerically at $x \approx 0.0785$, and the relative error is monotone in $x$
(\cref{fig:iv:mental-arithmetic:rule72}).
\end{proof}

\begin{omfigure}
\begin{tikzpicture}
\begin{axis}[width=10cm,height=5.2cm,xlabel={\small rate per period (\%)},ylabel={\small relative error (\%)},
  xmin=1,xmax=25,ymin=-12,ymax=5,xtick={1,5,10,15,20,25},
  tick label style={font=\footnotesize},grid=major,grid style={gray!20},table/col sep=comma,
  legend cell align=left,legend style={font=\scriptsize,at={(0.5,-0.32)},anchor=north,
  /tikz/every even column/.append style={column sep=8pt}},legend columns=3]
\addplot[omDef,thick] table[x=rate,y=r72]{figdata/interviews/08-mental-arithmetic/rule72.csv};
\addplot[omProp,thick,dashed] table[x=rate,y=r70]{figdata/interviews/08-mental-arithmetic/rule72.csv};
\addplot[omThm,thick,densely dotted] table[x=rate,y=r693]{figdata/interviews/08-mental-arithmetic/rule72.csv};
\draw[gray] (axis cs:1,0) -- (axis cs:25,0);
\legend{rule of 72,rule of 70,rule of 69.3}
\end{axis}
\end{tikzpicture}

\omcaption{Relative error of three rules for the doubling time against the exact $\ln 2/\ln(1+r)$, for rates
compounded once a period. The \omterm{def:iv:mental-arithmetic:rule72}{rule of 72} crosses zero near 7.85\%, the rule of 70 near 2\%; the rule of 69.3
is always below. Data: \texttt{fig\_iv\_rule72.py}.}\label{fig:iv:mental-arithmetic:rule72}
\end{omfigure}

\begin{example}[Compounding by logarithms]\label{ex:iv:mental-arithmetic:compound}
$1.03^{10} = e^{10\ln 1.03}$, and $\ln 1.03 \approx 0.03 - 0.00045 = 0.02955$, so the exponent is about 0.2955
and the result about $e^{0.3} e^{-0.0045} \approx 1.350 \times 0.9955 \approx 1.344$.
\end{example}

\section{Checking an answer}

A fast answer is useful only if it is checked as fast. Three checks cost a second each: the order of
magnitude (count the digits, or compare with a round product), the last digit (it is the last digit of the
product of the last digits), and the residue modulo 9.

\begin{definition}[Casting out nines]\label{def:iv:mental-arithmetic:nines}
\emph{Casting out nines}\index{casting out nines} checks an integer calculation by replacing each number with the
remainder of its digit sum on division by 9, which is the number's residue modulo 9, and checking that the
calculation holds for the residues.
\end{definition}

A product passes if the residues of the factors multiply to the residue of the result: $5\,824$ has digit sum
19, residue 1. The check fails to see errors that change the result by a multiple of 9, including the commonest
slip, the transposition of two digits (\cref{iq:iv:mental-arithmetic:9}): a passed check is evidence, a failed
check is proof.

\section{Worked answers}

\begin{example}[A P\&L in mixed units]\label{ex:iv:mental-arithmetic:pnl}
\emph{``You bought 1\,250 contracts; each tick is worth 12.50; the market moved three ticks in your favour. What did you
make?''} Both factors are eighths of powers of ten: $1\,250 = 10\,000/8$ and $12.5 = 100/8$, so $1\,250 \times 12.5 =
10^6/64 = 15\,625$, and three ticks make $46\,875$. \emph{Check:} the answer lies between $1\,000 \times 10 \times 3 =
30\,000$ and $1\,300 \times 13 \times 3 = 50\,700$, and closer to the top because both factors were rounded down by more
than they were rounded up. Saying the bounds takes three seconds and catches a slipped zero.
\end{example}

\begin{example}[Fees on a fund's return]\label{ex:iv:mental-arithmetic:fees}
\emph{``A fund charges 2\% of starting assets and 20\% of the gain above that fee. It made 12\% before fees. What did the
investor make, and what gross return would net 10\%?''} On 100: the fee takes 2, leaving a gain of 10, of which 20\% is
2, so the investor nets 8\%. In general the net is $0.8(g - 2)$ for a gross gain of $g$ per 100, so netting 10 needs $g - 2
= 12.5$: a gross return of 14.5\%. \emph{Check:} at $g = 2$ the investor nets zero and the fund keeps everything, which is
what the formula says. Fee structures differ in the order of the two fees and in hurdles; state the one assumed.
\end{example}

\begin{example}[From a monthly to a yearly return]\label{ex:iv:mental-arithmetic:annualise}
\emph{``A strategy returns 1.5\% a month. What is that compounded over a year?''} Take logarithms:
$\ln 1.015 \approx 0.015 - 0.015^2/2 = 0.0148875$; twelve months give $0.17865$; then $e^{0.17865} \approx 1 + 0.17865 +
0.01596 + 0.00095 \approx 1.1956$, so about 19.6\% a year, against 18\% for the simple sum. The exact value is
$1.015^{12} = 1.19562$ to five decimals. \emph{Check:} compounding must add something to $12 \times 1.5\% = 18\%$, and the
excess, about $\binom{12}{2} \times 0.015^2 \approx 1.5\%$, matches the first cross term.
\end{example}

\section{Question bank}

\begin{interviewq}[$\star$ \iqroles{trader} \iqfirm{market maker}]\label{iq:iv:mental-arithmetic:1}
$47 \times 53$.
\end{interviewq}

\begin{interviewq}[$\star$ \iqroles{trader} \iqfirm{market maker}]\label{iq:iv:mental-arithmetic:2}
$96 \times 104$ and $98 \times 97$.
\end{interviewq}

\begin{interviewq}[$\star$ \iqroles{trader, risk} \iqfirm{proprietary firm}]\label{iq:iv:mental-arithmetic:3}
$65^2$, and without writing anything, $75^2 - 65^2$.
\end{interviewq}

\begin{interviewq}[$\star$ \iqroles{trader, bank} \iqfirm{bank}]\label{iq:iv:mental-arithmetic:4}
12.5\% of 360, then 36\% of 25.
\end{interviewq}

\begin{interviewq}[$\star$ \iqroles{trader, researcher} \iqfirm{any}]\label{iq:iv:mental-arithmetic:5}
$\tfrac37$ and $\tfrac5{13}$ as decimals to three places.
\end{interviewq}

\begin{interviewq}[$\star\star$ \iqroles{researcher, trader} \iqfirm{any}]\label{iq:iv:mental-arithmetic:6}
$\sqrt{50}$ to four decimal places, with a bound on your error.
\end{interviewq}

\begin{interviewq}[$\star\star$ \iqroles{trader, bank} \iqfirm{bank}]\label{iq:iv:mental-arithmetic:7}
By the \omterm{def:iv:mental-arithmetic:rule72}{rule of 72}, how long does money take to double at 1\%, 6\% and 24\% a year? Which of the three answers is
furthest from the exact value, and in which direction?
\end{interviewq}

\begin{interviewq}[$\star\star$ \iqroles{trader} \iqfirm{market maker}]\label{iq:iv:mental-arithmetic:8}
A strategy earns 3.2 basis points on a daily turnover of 250 million. What does it earn a day, and in a year of
252 trading days?
\end{interviewq}

\begin{interviewq}[$\star\star$ \iqroles{trader, developer} \iqfirm{proprietary firm}]\label{iq:iv:mental-arithmetic:9}
Check $3\,847 \times 29 = 111\,563$ by \omterm{def:iv:mental-arithmetic:nines}{casting out nines}. A colleague wrote $111\,653$; does it pass? What does
this tell you about the check?
\end{interviewq}

\begin{interviewq}[$\star\star$ \iqroles{trader, risk} \iqfirm{any}]\label{iq:iv:mental-arithmetic:10}
A price rises 20\% then falls 20\%. Where is it relative to its start? What fall, after a 25\% rise, brings it
back exactly?
\end{interviewq}

\begin{interviewq}[$\star\star\star$ \iqroles{trader, bank} \iqfirm{bank}]\label{iq:iv:mental-arithmetic:11}
$\ln 1.07$ to four decimal places without a calculator, and from it the exact doubling time at 7\% a year to two
decimals.
\end{interviewq}

\begin{interviewq}[$\star\star\star$ \iqroles{trader, researcher} \iqfirm{market maker}]\label{iq:iv:mental-arithmetic:12}
$1.03^{10}$ to three decimal places.
\end{interviewq}

\begin{interviewq}[$\star\star\star$ \iqroles{trader} \iqfirm{market maker}]\label{iq:iv:mental-arithmetic:13}
$0.0625 \times 4.8 \times 125$, in under ten seconds. Say the rewrite you used.
\end{interviewq}

\begin{interviewq}[$\star\star\star$ \iqroles{researcher, bank} \iqfirm{any}]\label{iq:iv:mental-arithmetic:14}
Explain why the \omterm{def:iv:mental-arithmetic:rule72}{rule of 72} uses 72 rather than 69.3, and estimate the rate at which it is exact. Which rule
would you use for continuously compounded rates?
\end{interviewq}

\begin{omsources}
\item The rules of the chapter are elementary; the error bound of the linearised square root is the Lagrange
remainder of the Taylor series.
\item One Quant Book~1, chapter~7 (carry and P\&L arithmetic); One Quant Book~2, chapter~3 (compounding
conventions).
\end{omsources}
